\paragraph{Claim verdicts.}
Table~\ref{tab:verdict} lists the gap between our results and the paper's for
each claim, together with our verdict.

\begin{table}[h]
\centering\small
\caption{Claim verdicts. Gaps are ours minus the paper's (absolute, relative).}\label{tab:verdict}
\begin{tabularx}{\linewidth}{lXrrrX}
\toprule
Claim & Metric & Paper & Ours & Gap & Verdict \\
\midrule
C1 & avg.\ slowdown R+W, 7 counterparts & \PaperMapRW & \OurMapRW$\pm$\OurMapRWsd & \GapRWMap{} (\GapRWMapRel) & \textbf{Supported} (avg.\ only) \\
   & avg.\ slowdown R+W, all programs & 1.73 & \OurAvgRW$\pm$\OurAvgRWsd & \GapRWAll{} (\GapRWAllRel) & (higher; see text) \\
   & perl / XML slowdown & 2.60 / 2.67 & \RWPerl{} / \RWLibxml & \GapPerl{} / \GapXml & perl yes, XML lower \\
C2 & avg.\ slowdown W, 7 counterparts & \PaperMapWO & \OurMapWO$\pm$\OurMapWOsd & \GapWOMap{} (\GapWOMapRel) & \textbf{Supported} \\
   & avg.\ slowdown W, all programs & 1.26 & \OurAvgWO$\pm$\OurAvgWOsd & \GapWOAll{} (\GapWOAllRel) & (higher; see text) \\
C3 & memory increase (total) & 3.37 & \OurMemTotal & \GapMem{} (\GapMemRel) & \textbf{Not supported} (magnitude) \\
   & stack $>$5\% larger for few programs & 4 of 19 & \StackNoticeable{} of 11 & -- & \textbf{Supported} \\
   & binary size (code) & 1.5--3.2, avg 2.5 & \CodeMin--\CodeMax, avg \CodeAvg & -- & \textbf{Supported} (slightly larger) \\
\bottomrule
\end{tabularx}
\end{table}

\paragraph{C1: the slowdown is reproduced for comparable programs.}
For the seven programs with a SPEC counterpart, our average slowdown
(\OurMapRW$\times$) is within \GapRWMapRel{} of the paper's value for the same
seven benchmarks (\PaperMapRW$\times$). The average over all 11 programs is
\GapRWAllRel{} higher than the paper's 73\%. This gap is a property of the
workload mix, not of the tool. Even in the paper, the seven SPEC programs we
chose average \PaperMapRW$\times$, well above its 19-program mean of
1.73$\times$: we picked the best-known programs, and they include the paper's
two worst cases, while the cheapest SPEC programs (libquantum 1.08$\times$,
omnetpp 1.40$\times$, lbm 1.42$\times$) have no counterpart in our set. Our
four other programs add two bytecode interpreters (\cmd{lua} \RWLua$\times$,
\cmd{sqlite} \RWSqlite$\times$) and a template-heavy C++ code (\cmd{tramp3d}
\RWTrampd$\times$), which are all slowed more than the paper's average.

The paper explains its two worst cases by malloc intensity and many small
memory accesses. We tested this explanation by counting heap allocations in
the base builds (Table~\ref{tab:alloc}). The rank correlation between
allocation rate and slowdown is $\rho=\AllocSpearman$. The two most slowed
programs (\cmd{perl}, \cmd{lua}) are both among the \NumHiAlloc{} programs
with more than $10^6$ allocations/s; \cmd{libxml2} is the exception, discussed
below.

\paragraph{Per-program agreement is weaker than the average.}
The match of the seven-program average partly results from errors that
cancel. \cmd{stockfish} (\RWStockfish$\times$ against 1.60$\times$ for sjeng)
and \cmd{x264} (\RWX$\times$ against 1.46$\times$ for h264ref) are much slower
than their counterparts. \cmd{bzip2}, \cmd{hmmer} and \cmd{libxml2} are faster
than theirs. Only \cmd{perl} and \cmd{gnugo} are within about 0.2 of the
paper. We think the large gaps come from how the programs are built, but we
did not test this:
\begin{itemize}
\item \cmd{x264} is compiled with \cmd{--disable-asm}, so its hand-written
SIMD kernels are replaced by C loops that \asan{} instruments;
\item \cmd{stockfish} spends its time in a vectorised (AVX2) neural-network
evaluation with dense array accesses, unlike the 2006 sjeng;
\item \cmd{libxml2} is a different program from Xalan, discussed next.
\end{itemize}
The paper's claim is about the SPEC average, and our like-for-like average
agrees with it. Its per-benchmark bars cannot be reproduced with substitute
programs.

The clearest miss is the paper's worst case, xalancbmk. Its counterpart
\cmd{libxml2} is slowed \RWLibxml$\times$, against 2.67$\times$. libxml2 is the
same kind of application but a different program: it is written in C, not
template-heavy C++, and probably allocates fewer objects per byte of input
than Xalan. The comparison is therefore domain-level only, and we do not count
it as a contradiction.

\begin{table}[h]
\centering\small
\caption{Allocation profile of the base builds (\cmd{harness/malloccount.c},
\cmd{results/raw/malloc-profile/}) and slowdown (Table~\ref{tab:time}).}\label{tab:alloc}
\begin{tabular}{lrrrr}
\toprule
Benchmark & Heap allocations & Allocs/s & $\le$64 B & Slowdown R+W \\
\midrule
\texttt{perl} & 8,696,040 & 3,014,386 & 92\% & 2.82 \\
\texttt{bzip2} & 30 & 10 & 40\% & 1.48 \\
\texttt{gnugo} & 7,289 & 2,911 & 0\% & 1.53 \\
\texttt{hmmer} & 12,518 & 3,271 & 13\% & 1.56 \\
\texttt{stockfish} & 7,491 & 3,218 & 51\% & 2.33 \\
\texttt{x264} & 222 & 71 & 3\% & 2.22 \\
\texttt{libxml2} & 11,514,737 & 4,663,675 & 15\% & 1.92 \\
\texttt{tramp3d} & 3,784,225 & 1,797,009 & 69\% & 2.14 \\
\texttt{lulesh} & 30,063 & 11,863 & 0\% & 1.56 \\
\texttt{lua} & 30,775,214 & 10,730,674 & 100\% & 2.47 \\
\texttt{sqlite} & 855,613 & 329,373 & 30\% & 1.98 \\
\bottomrule
\end{tabular}

\end{table}

\paragraph{C2: removing read checks helps less for allocation-heavy programs.}
For the seven counterparts, writes-only instrumentation gives \OurMapWO$\times$
against \PaperMapWO$\times$ in the paper. Over all programs, the writes-only
mode keeps \WOShare\% of the full overhead, against \PaperWOShare\% in the
paper. The share depends on allocation intensity:
\begin{itemize}
\item for the \NumHiAlloc{} programs with more than $10^6$ allocations/s, it
is \WOShareHi\%;
\item for the others, it is \WOShareLo\%.
\end{itemize}
This is expected. Disabling read checks does not remove the work that \asan{}
does on every \cmd{malloc} and \cmd{free}: a slower allocator, plus poisoning
and unpoisoning the redzones in shadow memory. The rank correlation between
allocation rate and writes-only slowdown is $\rho=\AllocSpearmanWO$, even
stronger than for C1.

\paragraph{C3: memory, a measurement-method gap.}
The paper's metric no longer works. Current \asan{} reserves its whole shadow
region (about 20\,TB of virtual address space, mapped with
\cmd{MAP\_NORESERVE}) at start-up, so VmPeak is \VmPeakAsanTB\,TB for every
instrumented process (Table~\ref{tab:mem}). We therefore compare peak resident
memory (VmHWM). It grows \OurMemTotal$\times$ in total, much less than the
paper's 3.37$\times$. Two pieces of evidence explain this:
\begin{enumerate}
\item \emph{Constant overheads dominate small programs.} The paper's largest
ratios come from benchmarks with tiny footprints (povray 17$\rightarrow$371\,MB,
hmmer 33$\rightarrow$582\,MB), which indicates an additive cost of hundreds of
MB per process in the 2012 implementation (Table~1 of the paper). Measured
the same way, our additive cost is a few MB (e.g.\ +\MemAddHmmer\,MB for
\cmd{hmmer}, +\MemAddBzip\,MB for \cmd{bzip2}). Our ratio is largest for the
smallest program (\cmd{hmmer}, \MemHmmer$\times$) and smallest where large
buffers dominate (\cmd{stockfish} \MemStockfish$\times$, \cmd{x264}
\MemX$\times$). The pattern matches the paper's; only the constant has
shrunk.
\item \emph{Run-time settings matter more than instrumentation.} With
today's default options (256\,MB quarantine, 30-frame allocation stack traces,
16-byte minimum redzone), the total increase is \CtrlMemDefault$\times$, and
up to \CtrlMemDefMax$\times$ for \cmd{\CtrlMemDefMaxName}
(Table~\ref{tab:controls}). The paper excludes the quarantine (``constant-size
overhead \dots{} which we did not count''), and so do we, using
\cmd{quarantine\_size\_mb=0}.
\end{enumerate}
Whether redzones are the main \emph{variable} part of the overhead is tested
in \S\ref{sec:ext}.

The stack result agrees with the paper. Only \cmd{gnugo} grows noticeably
(1.30$\times$, against 1.43$\times$ for gobmk). Our \cmd{perl} shows no change,
whereas perlbench grows 3.06$\times$; the most likely reason is that our
workload recurses much less than SPEC's mail-processing inputs. Code growth,
once the constant \RuntimeKB\,KB run-time is excluded, matches the paper's
range.

\paragraph{Controlled re-runs: what changed since 2012?}
Table~\ref{tab:controls} varies one factor at a time, relative to the
\cmd{asan} variant (\CtrlRuns{} runs).
\begin{itemize}
\item \emph{Today's default run-time options} raise the average slowdown by
\CtrlDefDelta. Most of the increase falls on allocation-heavy programs
(\cmd{perl} \CtrlPapPerl$\rightarrow$\CtrlDefPerl$\times$, \cmd{lua}
\CtrlPapLua$\rightarrow$\CtrlDefLua$\times$), because collecting a stack trace
at every \cmd{malloc}/\cmd{free} is expensive. A reproduction that leaves the
options at their defaults would therefore report a larger slowdown than the
paper for reasons unrelated to the instrumentation.
\item \emph{Removing post-2012 instrumentation} (\cmd{asanlean}: no
use-after-scope poisoning, no stack-use-after-return hooks) lowers the average
by only \CtrlLeanDelta, which is within run-to-run noise. Twelve years of
compiler changes have thus left the cost of the core instrumentation roughly
where the paper measured it.
\end{itemize}

\begin{table}[h]
\centering\small
\caption{Controlled re-runs: slowdown (mean $\pm$ s.d., \CtrlRuns{} runs) and
memory increase (total peak RSS). ``paper opts'' = \cmd{asan} build with the
paper's options; ``default opts'' = same binary with today's defaults; ``lean
build'' = without use-after-scope and stack-use-after-return
instrumentation. Raw data: \cmd{results/raw/controls/}.}\label{tab:controls}
\begin{tabular}{lrrrrr}
\toprule
Benchmark & paper opts & default opts & lean build & mem paper & mem default \\
\midrule
\texttt{perl} & 2.60$\pm$0.57 & 3.87$\pm$0.85 & 2.64$\pm$0.47 & 1.23 & 5.37 \\
\texttt{bzip2} & 1.60$\pm$0.17 & 1.55$\pm$0.19 & 1.56$\pm$0.15 & 1.47 & 1.77 \\
\texttt{gnugo} & 1.80$\pm$0.30 & 2.06$\pm$0.29 & 1.67$\pm$0.31 & 1.63 & 4.29 \\
\texttt{hmmer} & 1.74$\pm$0.43 & 1.64$\pm$0.27 & 1.62$\pm$0.21 & 2.43 & 5.05 \\
\texttt{stockfish} & 2.08$\pm$0.56 & 2.37$\pm$0.60 & 1.92$\pm$0.30 & 1.03 & 1.36 \\
\texttt{x264} & 2.65$\pm$0.51 & 2.31$\pm$0.23 & 2.18$\pm$0.27 & 1.34 & 1.76 \\
\texttt{libxml2} & 1.94$\pm$0.21 & 2.16$\pm$0.44 & 1.79$\pm$0.08 & 1.64 & 2.30 \\
\texttt{tramp3d} & 2.01$\pm$0.17 & 2.45$\pm$0.26 & 2.18$\pm$0.29 & 1.32 & 21.28 \\
\texttt{lulesh} & 1.49$\pm$0.21 & 1.59$\pm$0.36 & 1.40$\pm$0.26 & 1.40 & 13.87 \\
\texttt{lua} & 2.55$\pm$0.43 & 3.60$\pm$0.35 & 2.54$\pm$0.31 & 1.92 & 13.62 \\
\texttt{sqlite} & 2.19$\pm$0.25 & 2.52$\pm$0.08 & 1.96$\pm$0.21 & 1.56 & 7.10 \\
\midrule
Average & 2.06$\pm$0.15 & 2.37$\pm$0.15 & 1.95$\pm$0.07 & 1.47 & 3.96 \\
\bottomrule
\end{tabular}

\end{table}

\paragraph{Measurement noise.}
Our machine is a cloud VM, and the hypervisor regularly runs other tenants on
our vCPU (``steal'' time in \cmd{/proc/stat}). The median coefficient of
variation of the base run times is \MedBaseCV\%, so per-program standard
deviations are large (up to \MaxRWsd). Two checks show that the averages are
robust to this noise:
\begin{itemize}
\item computing slowdowns from medians instead of means gives
\MedianAvgRW$\times$ / \MedianAvgWO$\times$;
\item in the \StealRuns{} repetitions where we recorded per-run steal time
(\StealPct\% of wall time on average, up to \StealMaxPct\%), subtracting
steal from wall time gives an average R+W slowdown of
\StealAdjRW$\pm$\StealAdjRWsd$\times$, against \StealRawRW$\times$
unadjusted, and lowers the median coefficient of variation of base run
times from \StealCVraw\% to \StealCVadj\%.
\end{itemize}

\paragraph{Limitations.}
\begin{itemize}
\item We did not run SPEC CPU2006. Agreement on ``counterparts'' shows that
the claims hold for similar programs, not that the paper's exact numbers are
reproduced.
\item Inputs are short (2--4\,s base). Start-up costs, which \asan{} pays on
every run (shadow reservation, allocator set-up), weigh more than with
minutes-long SPEC runs.
\item We used one 2-vCPU VM with noisy neighbours and one compiler version,
on x86-64 only.
\item The 2012 tool could not be rebuilt, so we cannot separate ``different
\asan{}'' from ``different workloads'' completely. The \cmd{asanlean} control
removes only the features we know about.
\item The paper's per-benchmark values are digitized ($\pm$0.03), and its
memory metric could not be applied.
\end{itemize}
